\section{An actual central-ample auxiliary line on the same normalized PL model}

\begin{theorem}[The original normalized tautological twist is ample]
\label{thm:pl-original-normalized-tautological-auxiliary-ample}
\lean{EnvelopeProperty.ProjSubvariety.normalizedSimultaneousTautologicalTwist_exists_ample}
\leanok
\uses{thm:pl-original-native-blowup-tautological-very-ample,thm:pl-auxiliary-native-finite-pullback-ample}
Let $I,J$ be the original homogeneous PL ideals and let the
original finite-log condition hold. On the same original
normalized simultaneous blowup, put
\[
 K=\nu^*((IJ)\mathcal O_{\operatorname{Bl}_{IJ}}),\qquad
 H_q=q_X^*\mathcal O_X(q)\otimes K.
\]
There is an internally chosen integer $q>0$ for which this
original whole $H_q$ is natively ample over $k[t]$.
\end{theorem}
\begin{proof}
\leanok
The original constant projective family has its actual closed
polynomial projective embedding and exact parameter-base triangle.
Apply the preceding native whole-blowup theorem to its literal
original product ideal $IJ$. It chooses $q>0$ and proves very
ampleness of the original tautological ideal twisted by the same
polynomial $\mathcal O(q)$ on the original all-affine blowup.

Its actual normalization morphism $\nu$ is finite and its
original composite structure map is projective. The finite
pullback theorem therefore makes the whole original pullback of
this same line ample. The canonical normalization Cartier-module
comparison identifies its ideal factor with the literal original
$K$. The already proved coefficient-twist comparison identifies
its polynomial twist with the same $q_X^*\mathcal O_X(q)$,
retaining the exponent $q$ exactly. Monoidal pullback combines
these whole isomorphisms into the original $H_q$. Transport
ampleness through this constructed isomorphism. No arbitrary
embedding line is declared equivariant, and no auxiliary ample
line or degree-match certificate has been supplied.
\end{proof}

\begin{theorem}[The same auxiliary line is ample on the original central fibre]
\label{thm:pl-original-normalized-tautological-central-ample}
\lean{EnvelopeProperty.ProjSubvariety.normalizedSimultaneousTautologicalTwist_exists_central_ample}
\leanok
\uses{thm:pl-original-normalized-tautological-auxiliary-ample,thm:pl-auxiliary-native-ampleness-base-change}
The same internally chosen $q>0$ gives native ampleness of the
same $H_q$ on the actual original central fibre, the canonical
pullback of the same model along $k[t]\to k$, $t\mapsto0$.
\end{theorem}
\begin{proof}
\leanok
Use the actual central Cartesian square and the preceding whole
coefficient base-change theorem. Its result is native ampleness
over $k$, of the original central inclusion's whole pullback of
the same $H_q$. The central fibre may be reducible or nonreduced.
\end{proof}

\begin{theorem}[Central ampleness of the auxiliary polarization on the normalized test configuration]
\label{thm:pl-original-normalized-tc-tautological-central-ample}
\lean{EnvelopeProperty.ProjSubvariety.normalizedSimultaneousNativeTestConfiguration_exists_central_ample_tautologicalTwist}
\leanok
\uses{thm:pl-original-normalized-tautological-central-ample,def:pl-original-normalized-tc-tautological-line-data}
Assume moreover that $I,J$ are invariant under $t\mapsto at$, contain powers of
$t$, and that $X$ is projectively normal, so that the normalized simultaneous
model is a test configuration $\mathcal C$ of $X$. Then for some $q>0$ the
linearized line $H_q$ (with fibre-one line $\mathcal O_X(q)$) has ample
restriction to the central fibre $\mathcal C_0$.
\end{theorem}
\begin{proof}
\leanok
The projection of $\mathcal C$ is the structure morphism of the normalized
simultaneous model, so this is the preceding theorem.
\end{proof}
